% Generated draft: requires mathematical and source review.
\documentclass[11pt]{article}
\usepackage[margin=25mm]{geometry}
\usepackage{fontspec}
\setmainfont{Latin Modern Roman}
\usepackage{amsmath,amssymb,mathtools,unicode-math}
\setmathfont{Latin Modern Math}
\usepackage{enumitem,needspace,xurl,hyperref}
\setlength{\parindent}{0pt}
\setlength{\parskip}{5pt}
\setlength{\emergencystretch}{3em}
\newcommand{\N}{\mathbb{N}}
\newcommand{\Z}{\mathbb{Z}}
\newcommand{\Q}{\mathbb{Q}}
\newcommand{\R}{\mathbb{R}}
\newcommand{\C}{\mathbb{C}}
\newcommand{\F}{\mathbb{F}}
\newcommand{\A}{\mathbb{A}}
\newcommand{\PP}{\mathbb{P}}
\newcommand{\E}{\mathbb{E}}
\newcommand{\eps}{\varepsilon}
\newcommand{\dd}{\,\mathrm{d}}
\begin{document}
\section*{UnsolvedMath ID: 30006118}
\textbf{Source problem code:} catalogue entry
\subsection{Definitions and mathematical statement}
Retain the notation of the supplied report extract: \(G\) is the group under consideration, \(\Lambda\) is the coefficient object, and \(\operatorname{Bun}_G\), \(\operatorname{Isoc}_G\), and \(\operatorname{Par}_G\) denote the objects occurring there. Write
\[
\mathcal C_{G,\Lambda}=\operatorname{Shv}(\operatorname{Bun}_G,\Lambda),\qquad
\mathcal D_{G,\Lambda}=\operatorname{Shv}(\operatorname{Isoc}_G,\Lambda),\qquad
\mathcal P_G=\operatorname{Perf}(\operatorname{Par}_G).
\]
Here \(\operatorname{Perf}\) denotes the category of perfect complexes on the intended parameter object. The supplied extract does not define that object, the two moduli objects, the sheaf categories, their categorical enhancements, or their coefficient conventions. It also does not specify the admissible groups, base fields, or coefficients. These are unresolved hypotheses, not permission to quantify over arbitrary groups or coefficient rings.

The target is Conjecture 3 of the supplied extract: for every instance allowed by the intended report hypotheses, construct a canonical equivalence
\[
\Phi_{G,\Lambda}\colon
\operatorname{Shv}(\operatorname{Bun}_G,\Lambda)
\xrightarrow{\ \sim\ }
\operatorname{Shv}(\operatorname{Isoc}_G,\Lambda)
\]
that is \(\operatorname{Perf}(\operatorname{Par}_G)\)-linear. In particular, for the intended actions of \(\mathcal P_G\), this requires coherent natural identifications
\[
\Phi_{G,\Lambda}(P\cdot A)\simeq P\cdot\Phi_{G,\Lambda}(A)
\qquad(P\in\mathcal P_G,\ A\in\mathcal C_{G,\Lambda}),
\]
compatible with the unit and tensor product. The actions and the intended meaning of canonical are not fully specified in the supplied record. Consequently, this is a conditional transcription of the conjectural target; a fully quantified formulation requires those missing definitions and hypotheses.

The surrounding passage supplies the following qualifications and construction context. It describes the full equivalence as conjectural, while attributing to its reference [17] the theorem that an unspecified morphism (1), constructed from the spectral action of Fargues--Scholze, agrees with the morphism of Genestier--Lafforgue. Neither morphism (1) nor the bibliographic identity of [17] is supplied, so this comparison cannot be restated precisely here.

The passage further attributes to spectral decomposition theorems [6, 1] an equivalence between the data of the spectral action and that of Hecke functors. In its notation, functorially in finite sets \(I\), these are tensor functors
\[
T_I\colon
\operatorname{Rep}_{\Lambda}({}^{L}G)^{\otimes I}
\longrightarrow
\operatorname{End}(\operatorname{Shv}(\operatorname{Isoc}_G))
\otimes\operatorname{Lisse}(D^{\times})^{\otimes I},
\qquad D^{\times}=\operatorname{Spec}(F((t))),
\]
where \(\operatorname{Lisse}(D^{\times})\) is described as the category of Weil local systems. The field \(F\), the precise meaning of \({}^{L}G\), and the tensor products and categorical conventions in this formula are not specified in the extract.

For \(I=\{\ast\}\), the passage introduces a deformation
\[
f\colon\operatorname{Isoc}_{G,\mathbb A^1}\longrightarrow\mathbb A^1
\]
with fiber over \(0\) equal to \(\operatorname{Isoc}_G\), and describes its fiber over \(\mathbb A^1\setminus\{0\}\) as \(\operatorname{Isoc}_G\times\operatorname{Hk}_{\mathrm{loc}}\), where \(\operatorname{Hk}_{\mathrm{loc}}\) is the local Hecke stack. The precise family or trivialization intended by this latter description is not supplied. It gives the construction formula
\[
T_{\{\ast\}}(V)(A)
=\Psi_f(A\otimes\operatorname{Sat}_V),
\qquad
V\in\operatorname{Rep}_{\Lambda}({}^{L}G),\quad
A\in\operatorname{Shv}(\operatorname{Isoc}_G),
\]
where \(\Psi_f\) is the nearby-cycle functor for \(f\), and \(\operatorname{Sat}_V\in\operatorname{Shv}(\operatorname{Hk}_{\mathrm{loc}})\) is the Satake sheaf. Pullbacks, tensor-product conventions, and coefficient hypotheses needed to interpret this formula are not specified.

For \(|I|>1\), the passage invokes nearby cycles over general bases, attributing the theory to Deligne, Laumon [16], Gabber, Orgogozo [18], and Illusie [12], with a reinterpretation by Hansen--Scholze [10]. It explains that control generally requires modifying the base, whereas its claimed canonical construction uses naive nearby cycles without a modification. It reports proving the required good behavior through contractibility, in an unspecified suitable sense, of a space controlling the difference between naive and non-naive nearby cycles. It also identifies the need to pass from the a priori inertia action \(I_F\subset W_F\) to the Weil-group action by coherently organizing partial Frobenius correspondences with nearby cycles. These are supplied contextual claims and qualifications; the extract does not provide enough definitions to turn them into separate precise assertions.
\subsection{Short English statement}
Construct the canonical equivalence between the sheaf categories on \(\operatorname{Bun}_G\) and \(\operatorname{Isoc}_G\), compatibly with their \(\operatorname{Perf}(\operatorname{Par}_G)\)-actions, throughout the intended range of groups and coefficients. That range and several essential definitions are missing from the extract. Supplied curation dated 2026-08-22 labels the problem partially solved and reports a torsion-coefficient comparison; this assessment has not been independently verified here.
\subsection{Sources}
\begin{itemize}
\item Supplied local UnsolvedMath registry record, permanent ID \(30006118\), title \emph{Sheaves on Bundles and Isocrystals}, exact source problem code \texttt{catalogue entry}; collection: Oberwolfach Reports Open Problems. The record gives proposed year 2024 and no named proposer.
\item Supplied report citation: \emph{Representations of \(p\)-adic Groups (2025)}, Oberwolfach Reports, pp.~3103--3178, DOI \texttt{10.4171/owr/2024/54}. \url{https://doi.org/10.4171/owr/2024/54}. The extracted target is labelled Conjecture 3, equation (2). The citation's 2025 title date and its 2024 report identifier are preserved without reconciliation.
\item Supplied literature source, labelled ``OWR report'': \url{https://publications.mfo.de/bitstream/handle/mfo/4270/OWR_2024_54.pdf?sequence=1}.
\item Supplied literature source, labelled ``Schematic and analytic LLC'': \url{https://arxiv.org/abs/2606.02799}. No authors or theorem number for the asserted comparison are supplied.
\item Literature assessment checked at \texttt{2026-08-22}: the supplied curation labels the problem \texttt{partially\_solved} and states that a 2026 primary preprint proves the comparison for torsion coefficients, while the full coefficient-independent equivalence is not established in complete generality. This is supplied dated curation, not an independently verified literature assessment; the record explicitly requests independent verification of time-sensitive status and cited claims.
\item The record labels the statement \texttt{corrected\_verified} and says it was reconstructed and checked against the official Oberwolfach report. This is supplied provenance, not verification performed for this draft.
\item Attribution retained from the passage: Fargues--Scholze and Genestier--Lafforgue for the compared constructions; spectral decomposition references [6, 1]; comparison reference [17]; Deligne, Laumon [16], Gabber, Orgogozo [18], Illusie [12], and Hansen--Scholze [10] for the nearby-cycle context. The corresponding bibliography entries are not included in the supplied record.
\end{itemize}
\end{document}
